% abstract.tex — Abstracts (IT + EN, ≤200 words each, Vademecum Rev06)
% Paste the final text into the \TODO markers below.

\chapter*{Sommario}
\addcontentsline{toc}{chapter}{Sommario}

eBPF (extended Berkeley Packet Filter) è diventato un componente fondamentale
dei moderni sistemi Linux, consentendo l'esecuzione sicura e sandboxed di programmi
utente all'interno del kernel per applicazioni di rete, osservabilità e sicurezza.
Il verifier eBPF~--- il componente del kernel responsabile della garanzia di
sicurezza dei programmi~--- rappresenta pertanto un obiettivo di alto valore:
falle logiche nel verifier possono permettere a codice non privilegiato di uscire
dalla sandbox e compromettere il sistema host.

Questa tesi presenta una pipeline automatizzata per il fuzzing coverage-guided del
verifier eBPF di Linux tramite un large language model (LLM) fine-tunato.
A partire da circa 2 milioni di programmi eBPF raccolti al confine FFI del kernel
mediante una versione modificata del fuzzer Buzzer, curiamo un dataset bilanciato
di 27\,000 campioni distribuiti tra le classi di errore del verifier ed eseguiamo
il fine-tuning di Qwen2.5-Coder-1.5B tramite QLoRA per generare assembly in
formato verifier-log sintatticamente valido.
Applichiamo poi l'apprendimento per rinforzo GRPO con un segnale di reward derivato
dalla copertura del codice kernel (KCOV), addestrando il modello a esplorare
percorsi del verifier più profondi e meno visitati.
Una formulazione del reward verdict-blind e basata sulla profondità evita il
problema di collasso del reward osservato quando programmi accettati e rifiutati
vengono premiati in modo diverso.

I risultati sperimentali mostrano che il modello fine-tunato migliora il pass-rate
rispetto alla baseline zero-shot, e che la fase di training RL coverage-guided
guida con successo l'esplorazione di nuovi percorsi del codice kernel.
Il sistema dimostra che il fuzzing guidato da LLM è un approccio valido per la
ricerca sulla sicurezza del verifier, e fornisce una pipeline open-source
riproducibile per lavori futuri.

\bigskip
\noindent\textbf{Parole chiave:} eBPF, fuzzing, modelli linguistici di grandi dimensioni,
apprendimento per rinforzo, copertura del codice kernel, KCOV, GRPO.

\cleardoublepage

\chapter*{Abstract}
\addcontentsline{toc}{chapter}{Abstract}

eBPF (extended Berkeley Packet Filter) has become a cornerstone of modern Linux
systems, enabling safe, sandboxed execution of user-supplied programs inside the
kernel for networking, observability, and security applications.
The eBPF verifier~--- the kernel component responsible for guaranteeing program
safety~--- is therefore a high-value target: logical flaws in the verifier can
allow unprivileged code to escape the sandbox and compromise the host system.

This thesis presents an automated pipeline for coverage-guided fuzzing of the Linux
eBPF verifier using a fine-tuned large language model (LLM).
Starting from approximately 2 million eBPF programs collected at the kernel FFI
boundary using a modified version of the Buzzer fuzzer, we curate a balanced
dataset of 27,000 samples across verifier error classes and fine-tune
Qwen2.5-Coder-1.5B via QLoRA to generate syntactically valid verifier-log-format
assembly.
We then apply GRPO reinforcement learning with a reward signal derived from kernel
code coverage (KCOV), training the model to explore deeper and less-visited
verifier paths.
A verdict-blind, depth-based reward formulation avoids the reward-collapse problem
observed when accepted and rejected programs are rewarded differently.

Experimental results show that the fine-tuned model improves pass-rate over the
zero-shot baseline, and that the coverage-guided RL training phase successfully
drives exploration of new kernel code paths.
The system demonstrates that LLM-guided fuzzing is a viable approach for verifier
security research, and provides a reproducible open-source pipeline for future work.

\bigskip
\noindent\textbf{Keywords:} eBPF, fuzzing, large language models, reinforcement learning,
kernel coverage, KCOV, GRPO.
